\section{OLS 斜率是怎么定出来的？}\label{sec:ols-slope}

\subsection{两个函数不能混用}\label{sec:prf-srf}

\StudyTerm{\PriorityMust}{总体回归函数}{$E(y\mid x)=\beta_0+\beta_1x$ 描述给定 $x$ 时 $y$ 的条件均值。}

\StudyTerm{\PriorityMust}{样本回归函数}{$\hat y=\hat\beta_0+\hat\beta_1x$ 是它的样本对应物，拟合的是手头这组数据。}

$\beta_1$ 是总体的未知参数，$\hat\beta_1$ 是用样本计算的估计量。总体不变时，重新抽样会改变估计量的取值，参数本身保持不变。

\subsection{一阶条件与斜率}\label{sec:foc}

最小二乘准则把 $\hat\beta_0,\hat\beta_1$ 取为使残差平方和
$SSR=\sum_i (y_i-\hat\beta_0-\hat\beta_1x_i)^2$ 最小的一对值。

记 $\hat S_{xy}=\sum_i (x_i-\bar x)(y_i-\bar y)$、$\hat S_{xx}=\sum_i (x_i-\bar x)^2$。
分别对截距和斜率求偏导并令其为零，得到
\begin{align*}
  \sum_i(y_i-\hat\beta_0-\hat\beta_1x_i)&=0,\\
  \sum_ix_i(y_i-\hat\beta_0-\hat\beta_1x_i)&=0.
\end{align*}
第一式给出 $\hat\beta_0=\bar y-\hat\beta_1\bar x$。从第二式减去第一式的 $\bar x$ 倍，再代入截距，得
$\sum_i(x_i-\bar x)[(y_i-\bar y)-\hat\beta_1(x_i-\bar x)]=0$，即

\begin{equation}
  \hat S_{xy}-\hat\beta_1\hat S_{xx}=0.
\end{equation}

两边解出斜率即得

\begin{equation}
  \hat\beta_1=\frac{\hat S_{xy}}{\hat S_{xx}}.\label{eq:ols-slope}
\end{equation}

整理时除以 $\hat S_{xx}$ 这一步骤要求 $\hat S_{xx}>0$，也就是 $x$ 在样本中不是常数。
若 $x$ 取值全部相同，截距可补偿任意斜率，无法分别确定两者。

截距的一阶条件同时给出 $\sum_i\hat u_i=0$，且 $\hat\beta_0+\hat\beta_1\bar x=\bar y$，因此含截距的 OLS 回归线过样本均值点。这些代数性质不需要误差项的分布假设。
